\section{Class-specific correlation estimates}
\label{sec:class-details}
Table~\ref{tab:class-correlations} reports the split-averaged correlations used to calculate effective support and their training-patient sensitivity intervals. These are descriptive class-level comparisons; no adjustment for multiple class comparisons is applied. Table~\ref{tab:raw-correlations} gives the corresponding raw estimates for each split. All estimates, including the resampled values, lie within $[0,1]$, so raw and clipped point estimates coincide.

\begin{table}[htbp]
\centering
\footnotesize
\setlength{\tabcolsep}{4pt}
\renewcommand{\arraystretch}{1.08}
\begin{tabularx}{\textwidth}{@{}Yrr@{}}
\toprule
Cancer type & Single direction & Full features \\
\midrule
Adrenocortical carcinoma & 0.500 [0.315, 0.620] & 0.433 [0.381, 0.473] \\
Bladder Urothelial Carcinoma & 0.508 [0.462, 0.549] & 0.521 [0.504, 0.537] \\
Brain Lower Grade Glioma & 0.429 [0.376, 0.477] & 0.484 [0.465, 0.498] \\
Breast invasive carcinoma & 0.448 [0.402, 0.489] & 0.511 [0.497, 0.523] \\
Cervical squamous cell / endocervical adenocarcinoma & 0.515 [0.456, 0.566] & 0.516 [0.495, 0.533] \\
Colon adenocarcinoma & 0.394 [0.310, 0.475] & 0.425 [0.395, 0.450] \\
Esophageal carcinoma & 0.410 [0.328, 0.481] & 0.518 [0.490, 0.538] \\
Glioblastoma multiforme & 0.337 [0.277, 0.393] & 0.400 [0.380, 0.418] \\
Head and neck squamous cell carcinoma & 0.329 [0.283, 0.373] & 0.383 [0.364, 0.400] \\
Kidney Chromophobe & 0.489 [0.365, 0.587] & 0.553 [0.512, 0.584] \\
Kidney renal clear cell carcinoma & 0.641 [0.581, 0.689] & 0.548 [0.522, 0.571] \\
Kidney renal papillary cell carcinoma & 0.564 [0.503, 0.612] & 0.555 [0.535, 0.571] \\
Liver hepatocellular carcinoma & 0.521 [0.468, 0.569] & 0.532 [0.513, 0.549] \\
Lung adenocarcinoma & 0.398 [0.360, 0.433] & 0.442 [0.426, 0.457] \\
Lung squamous cell carcinoma & 0.314 [0.275, 0.350] & 0.367 [0.353, 0.378] \\
Mesothelioma & 0.414 [0.279, 0.519] & 0.506 [0.457, 0.539] \\
Ovarian serous cystadenocarcinoma & 0.513 [0.404, 0.599] & 0.467 [0.428, 0.495] \\
Pancreatic adenocarcinoma & 0.328 [0.269, 0.383] & 0.434 [0.398, 0.464] \\
Pheochromocytoma and Paraganglioma & 0.636 [0.491, 0.721] & 0.582 [0.523, 0.619] \\
Prostate adenocarcinoma & 0.527 [0.441, 0.600] & 0.506 [0.476, 0.534] \\
Rectum adenocarcinoma & 0.530 [0.393, 0.626] & 0.520 [0.471, 0.560] \\
Sarcoma & 0.493 [0.419, 0.552] & 0.526 [0.500, 0.547] \\
Skin Cutaneous Melanoma & 0.462 [0.410, 0.511] & 0.542 [0.525, 0.555] \\
Stomach adenocarcinoma & 0.452 [0.404, 0.494] & 0.487 [0.468, 0.502] \\
Testicular Germ Cell Tumors & 0.493 [0.429, 0.553] & 0.500 [0.467, 0.530] \\
Thymoma & 0.358 [0.256, 0.444] & 0.477 [0.430, 0.512] \\
Thyroid carcinoma & 0.520 [0.468, 0.564] & 0.569 [0.551, 0.586] \\
Uterine Carcinosarcoma & 0.378 [0.218, 0.544] & 0.379 [0.330, 0.425] \\
Uterine Corpus Endometrial Carcinoma & 0.485 [0.439, 0.524] & 0.464 [0.448, 0.477] \\
Uveal Melanoma & 0.329 [0.252, 0.390] & 0.476 [0.440, 0.499] \\
\bottomrule
\end{tabularx}

\caption{Class-specific correlations averaged over three training splits, with separate 95\% training-patient sensitivity intervals for each measure. Full-feature correlation is higher in 22 of 30 classes; increases are not uniform across cancer types.}
\label{tab:class-correlations}
\end{table}

\clearpage
\begin{table}[htbp]
\centering
\footnotesize
\setlength{\tabcolsep}{4pt}
\renewcommand{\arraystretch}{1.08}
\begin{tabularx}{\textwidth}{@{}Yrrrrrr@{}}
\toprule
& \multicolumn{3}{c}{Single direction} & \multicolumn{3}{c}{Full features} \\
\cmidrule(lr){2-4}\cmidrule(l){5-7}
Cancer type & Split 0 & Split 1 & Split 2 & Split 0 & Split 1 & Split 2 \\
\midrule
Adrenocortical carcinoma & 0.501 & 0.473 & 0.527 & 0.431 & 0.427 & 0.441 \\
Bladder Urothelial Carcinoma & 0.503 & 0.515 & 0.507 & 0.529 & 0.524 & 0.510 \\
Brain Lower Grade Glioma & 0.448 & 0.435 & 0.404 & 0.488 & 0.479 & 0.485 \\
Breast invasive carcinoma & 0.446 & 0.444 & 0.454 & 0.511 & 0.513 & 0.510 \\
Cervical squamous cell / endocervical adenocarcinoma & 0.518 & 0.512 & 0.516 & 0.516 & 0.518 & 0.514 \\
Colon adenocarcinoma & 0.379 & 0.405 & 0.397 & 0.421 & 0.434 & 0.421 \\
Esophageal carcinoma & 0.438 & 0.405 & 0.387 & 0.529 & 0.521 & 0.503 \\
Glioblastoma multiforme & 0.350 & 0.293 & 0.367 & 0.400 & 0.391 & 0.410 \\
Head and neck squamous cell carcinoma & 0.340 & 0.310 & 0.337 & 0.387 & 0.377 & 0.387 \\
Kidney Chromophobe & 0.490 & 0.523 & 0.454 & 0.553 & 0.558 & 0.549 \\
Kidney renal clear cell carcinoma & 0.639 & 0.653 & 0.630 & 0.551 & 0.553 & 0.540 \\
Kidney renal papillary cell carcinoma & 0.569 & 0.572 & 0.551 & 0.553 & 0.559 & 0.553 \\
Liver hepatocellular carcinoma & 0.516 & 0.515 & 0.531 & 0.531 & 0.528 & 0.538 \\
Lung adenocarcinoma & 0.405 & 0.394 & 0.394 & 0.439 & 0.447 & 0.440 \\
Lung squamous cell carcinoma & 0.291 & 0.329 & 0.323 & 0.360 & 0.369 & 0.372 \\
Mesothelioma & 0.391 & 0.435 & 0.417 & 0.501 & 0.520 & 0.495 \\
Ovarian serous cystadenocarcinoma & 0.544 & 0.479 & 0.516 & 0.482 & 0.454 & 0.465 \\
Pancreatic adenocarcinoma & 0.338 & 0.300 & 0.345 & 0.436 & 0.429 & 0.439 \\
Pheochromocytoma and Paraganglioma & 0.603 & 0.651 & 0.654 & 0.565 & 0.588 & 0.594 \\
Prostate adenocarcinoma & 0.538 & 0.528 & 0.516 & 0.505 & 0.504 & 0.508 \\
Rectum adenocarcinoma & 0.551 & 0.501 & 0.538 & 0.522 & 0.513 & 0.526 \\
Sarcoma & 0.502 & 0.497 & 0.478 & 0.528 & 0.529 & 0.520 \\
Skin Cutaneous Melanoma & 0.467 & 0.466 & 0.452 & 0.544 & 0.541 & 0.540 \\
Stomach adenocarcinoma & 0.478 & 0.436 & 0.441 & 0.493 & 0.483 & 0.484 \\
Testicular Germ Cell Tumors & 0.507 & 0.508 & 0.465 & 0.516 & 0.503 & 0.482 \\
Thymoma & 0.352 & 0.365 & 0.356 & 0.476 & 0.476 & 0.478 \\
Thyroid carcinoma & 0.526 & 0.508 & 0.527 & 0.572 & 0.568 & 0.567 \\
Uterine Carcinosarcoma & 0.376 & 0.342 & 0.415 & 0.387 & 0.370 & 0.381 \\
Uterine Corpus Endometrial Carcinoma & 0.458 & 0.504 & 0.493 & 0.460 & 0.467 & 0.466 \\
Uveal Melanoma & 0.314 & 0.327 & 0.344 & 0.461 & 0.490 & 0.476 \\
\bottomrule
\end{tabularx}

\caption{Unclipped correlations for each class and training split, computed on the same sampled patients and patches for both measures. Every estimate is finite and positive. Split variation is descriptive because the training partitions overlap.}
\label{tab:raw-correlations}
\end{table}
